%====================================================%
% SEZIONE 1: Sistemi lineari di due equazioni
%====================================================%

\begin{question}[ID=1.1, subtitle={Sistemi lineari di due equazioni}, difficulty=facile, tags={sistemi, equazioni lineari}]
Risolvi i seguenti sistemi lineari di due equazioni:
\begin{multicols}{2}
\begin{enumerate}[leftmargin=3.2cm]
    \item
    \[
    \begin{cases}
    x+y=5\\
    x-y=1
    \end{cases}
    \]

    \item
    \[
    \begin{cases}
    2x+y=7\\
    x-y=2
    \end{cases}
    \]

    \item
    \[
    \begin{cases}
    x+2y=8\\
    x-y=2
    \end{cases}
    \]

    \item
    \[
    \begin{cases}
    3x+y=10\\
    x+y=6
    \end{cases}
    \]
\end{enumerate}
\end{multicols}
\begin{solution}
\begin{enumerate}
    \item $(x,y)=(3,2)$
    \item $(x,y)=(3,1)$
    \item $(x,y)=(4,2)$
    \item $(x,y)=(2,4)$
\end{enumerate}
\end{solution}
\end{question}

%----------------------------------------------------%

\begin{question}[ID=1.2, subtitle={Sistemi lineari con coefficienti interi}, difficulty=facile, tags={sistemi, equazioni lineari}]
Risolvi i seguenti sistemi lineari di due equazioni:
\begin{multicols}{2}
\begin{enumerate}[leftmargin=3.2cm]
    \item
    \[
    \begin{cases}
    2x+3y=12\\
    x-y=1
    \end{cases}
    \]

    \item
    \[
    \begin{cases}
    3x-2y=4\\
    x+y=3
    \end{cases}
    \]

    \item
    \[
    \begin{cases}
    4x+y=11\\
    2x-y=1
    \end{cases}
    \]

    \item
    \[
    \begin{cases}
    5x+2y=16\\
    x-2y=-4
    \end{cases}
    \]
\end{enumerate}
\end{multicols}
\begin{solution}
\begin{enumerate}
    \item $(x,y)=(3,2)$
    \item $(x,y)=(2,1)$
    \item $(x,y)=(2,3)$
    \item $(x,y)=(1,3)$
\end{enumerate}
\end{solution}
\end{question}

%----------------------------------------------------%

\begin{question}[ID=1.3, subtitle={Sistemi lineari con parentesi e frazioni}, difficulty=media, tags={sistemi, equazioni lineari}]
Risolvi i seguenti sistemi lineari di due equazioni:
\begin{multicols}{2}
\begin{enumerate}[leftmargin=3.2cm]
    \item
    \[
    \begin{cases}
    2(x+y)=10\\
    3x-y=5
    \end{cases}
    \]

    \item
    \[
    \begin{cases}
    \dfrac{x}{2}+\dfrac{y}{3}=3\\
    x-y=1
    \end{cases}
    \]

    \item
    \[
    \begin{cases}
    3(x-1)+2y=7\\
    x+y=3
    \end{cases}
    \]

    \item
    \[
    \begin{cases}
    \dfrac{2x-1}{3}+y=4\\
    x-\dfrac{y+1}{2}=2
    \end{cases}
    \]
\end{enumerate}
\end{multicols}
\begin{solution}
\begin{enumerate}
    \item $(x,y)=(3,2)$
    \item $(x,y)=\left(\tfrac{19}{5},\tfrac{14}{5}\right)$
    \item $(x,y)=(2,1)$
    \item $(x,y)=\left(\tfrac{13}{4},\tfrac54\right)$
\end{enumerate}
\end{solution}
\end{question}

%----------------------------------------------------%

\begin{question}[ID=1.4, subtitle={Sistemi lineari con parametro numerico e coefficienti frazionari}, difficulty=difficile, tags={sistemi, equazioni lineari}]
Risolvi i seguenti sistemi lineari di due equazioni:
\begin{multicols}{2}
\begin{enumerate}[leftmargin=3.2cm]
    \item
    \[
    \begin{cases}
    4x-3y=1\\
    5x+2y=19
    \end{cases}
    \]

    \item
    \[
    \begin{cases}
    \dfrac{3x+y}{2}=5\\
    \dfrac{x-y}{3}=1
    \end{cases}
    \]

    \item
    \[
    \begin{cases}
    2(x-2)-3(y+1)=4\\
    5x+2y=13
    \end{cases}
    \]

    \item
    \[
    \begin{cases}
    \dfrac{2x-3y}{4}=2\\
    \dfrac{x+y}{3}=2
    \end{cases}
    \]
\end{enumerate}
\end{multicols}
\begin{solution}
\begin{enumerate}
    \item $(x,y)=(3,1)$
    \item $(x,y)=\left(\tfrac{16}{5},\tfrac65\right)$
    \item $(x,y)=\left(\tfrac{37}{19},-\tfrac{4}{19}\right)$
    \item $(x,y)=\left(\tfrac{21}{5},\tfrac95\right)$
\end{enumerate}
\end{solution}
\end{question}

%----------------------------------------------------%

\begin{question}[ID=1.5, subtitle={Sistemi lineari: casi particolari}, difficulty=difficile, tags={sistemi, equazioni lineari}]
Risolvi e classifica i seguenti sistemi lineari di due equazioni:
\begin{multicols}{2}
\begin{enumerate}[leftmargin=3.2cm]
    \item
    \[
    \begin{cases}
    2x+4y=6\\
    x+2y=3
    \end{cases}
    \]

    \item
    \[
    \begin{cases}
    3x-6y=9\\
    x-2y=1
    \end{cases}
    \]

    \item
    \[
    \begin{cases}
    4x+2y=8\\
    2x+y=5
    \end{cases}
    \]

    \item
    \[
    \begin{cases}
    2x-3y=7\\
    4x-6y=14
    \end{cases}
    \]
\end{enumerate}
\end{multicols}
\begin{solution}
\begin{enumerate}
    \item Sistema indeterminato: infinite soluzioni.
    \item Sistema impossibile: nessuna soluzione.
    \item Sistema impossibile: nessuna soluzione.
    \item Sistema indeterminato: infinite soluzioni.
\end{enumerate}
\end{solution}
\end{question}


%====================================================%
% SEZIONE 2: Sistemi lineari con coefficienti frazionari
%====================================================%

\begin{question}[ID=2.1, subtitle={Sistemi con coefficienti frazionari}, difficulty=facile, tags={sistemi, equazioni lineari, frazioni}]
Risolvi i seguenti sistemi lineari di due equazioni:
\begin{multicols}{2}
\begin{enumerate}[leftmargin=3.2cm]
    \item
    \[
    \begin{cases}
    \dfrac{1}{2}x+\dfrac{1}{3}y=3\\
    x-y=1
    \end{cases}
    \]

    \item
    \[
    \begin{cases}
    \dfrac{1}{3}x+\dfrac{1}{2}y=3\\
    2x-y=2
    \end{cases}
    \]

    \item
    \[
    \begin{cases}
    \dfrac{1}{4}x+y=4\\
    \dfrac{1}{2}x-\dfrac{1}{3}y=1
    \end{cases}
    \]

    \item
    \[
    \begin{cases}
    \dfrac{1}{2}x+\dfrac{1}{5}y=4\\
    x-\dfrac{1}{2}y=-1
    \end{cases}
    \]
\end{enumerate}
\end{multicols}
\begin{solution}
\begin{enumerate}
    \item $(x,y)=(4,3)$
    \item $(x,y)=(3,4)$
    \item $(x,y)=(4,3)$
    \item $(x,y)=(4,10)$
\end{enumerate}
\end{solution}
\end{question}

%----------------------------------------------------%

\begin{question}[ID=2.2, subtitle={Sistemi con frazioni e numeratori letterali}, difficulty=media, tags={sistemi, equazioni lineari, frazioni}]
Risolvi i seguenti sistemi lineari di due equazioni:
\begin{multicols}{2}
\begin{enumerate}[leftmargin=3.2cm]
    \item
    \[
    \begin{cases}
    \dfrac{x+1}{2}+\dfrac{1}{3}y=3\\
    x-\dfrac{1}{3}y=2
    \end{cases}
    \]

    \item
    \[
    \begin{cases}
    \dfrac{1}{4}x+\dfrac{y-1}{3}=2\\
    x-\dfrac{1}{2}y=2
    \end{cases}
    \]

    \item
    \[
    \begin{cases}
    \dfrac{x-2}{3}+\dfrac{1}{2}y=3\\
    2x-y=6
    \end{cases}
    \]

    \item
    \[
    \begin{cases}
    \dfrac{1}{2}x+\dfrac{y+1}{4}=3\\
    \dfrac{x-1}{3}+y=4
    \end{cases}
    \]
\end{enumerate}
\end{multicols}
\begin{solution}
\begin{enumerate}
    \item $(x,y)=(3,3)$
    \item $(x,y)=(4,4)$
    \item $(x,y)=(5,4)$
    \item $(x,y)=(4,3)$
\end{enumerate}
\end{solution}
\end{question}

%----------------------------------------------------%

\begin{question}[ID=2.3, subtitle={Sistemi con frazioni in entrambe le equazioni}, difficulty=media, tags={sistemi, equazioni lineari, frazioni}]
Risolvi i seguenti sistemi lineari di due equazioni:
\begin{multicols}{2}
\begin{enumerate}[leftmargin=3.2cm]
    \item
    \[
    \begin{cases}
    \dfrac{2x-1}{3}+\dfrac{y+1}{2}=3\\
    \dfrac{1}{2}x-\dfrac{1}{3}y=0
    \end{cases}
    \]

    \item
    \[
    \begin{cases}
    \dfrac{x+y}{2}+\dfrac{1}{3}x=4\\
    x-\dfrac{1}{3}y=2
    \end{cases}
    \]

    \item
    \[
    \begin{cases}
    \dfrac{3x-y}{4}+\dfrac{1}{2}x=3\\
    \dfrac{1}{2}x-\dfrac{1}{4}y=0
    \end{cases}
    \]

    \item
    \[
    \begin{cases}
    \dfrac{1}{2}x-\dfrac{y-1}{3}=1\\
    \dfrac{x+2}{3}-\dfrac{1}{4}y=1
    \end{cases}
    \]
\end{enumerate}
\end{multicols}
\begin{solution}
\begin{enumerate}
    \item $(x,y)=(2,3)$
    \item $(x,y)=(3,3)$
    \item $(x,y)=(4,8)$
    \item $(x,y)=(4,4)$
\end{enumerate}
\end{solution}
\end{question}

%----------------------------------------------------%

\begin{question}[ID=2.4, subtitle={Sistemi con denominatori diversi}, difficulty=difficile, tags={sistemi, equazioni lineari, frazioni}]
Risolvi i seguenti sistemi lineari di due equazioni:
\begin{multicols}{2}
\begin{enumerate}[leftmargin=3.2cm]
    \item
    \[
    \begin{cases}
    \dfrac{1}{2}x+\dfrac{1}{3}y=5\\
    \dfrac{1}{4}x+\dfrac{1}{3}y=4
    \end{cases}
    \]

    \item
    \[
    \begin{cases}
    \dfrac{2}{3}x-\dfrac{1}{2}y=1\\
    \dfrac{1}{2}x+\dfrac{1}{3}y=5
    \end{cases}
    \]

    \item
    \[
    \begin{cases}
    \dfrac{3}{5}x+\dfrac{1}{2}y=4\\
    \dfrac{1}{5}x-\dfrac{1}{2}y=0
    \end{cases}
    \]

    \item
    \[
    \begin{cases}
    \dfrac{1}{3}x-\dfrac{2}{5}y=0\\
    \dfrac{1}{2}x+\dfrac{1}{5}y=4
    \end{cases}
    \]
\end{enumerate}
\end{multicols}
\begin{solution}
\begin{enumerate}
    \item $(x,y)=(4,9)$
    \item $(x,y)=(6,6)$
    \item $(x,y)=(5,2)$
    \item $(x,y)=(6,5)$
\end{enumerate}
\end{solution}
\end{question}

%----------------------------------------------------%

\begin{question}[ID=2.5, subtitle={Sistemi con frazioni complesse}, difficulty=difficile, tags={sistemi, equazioni lineari, frazioni}]
Risolvi i seguenti sistemi lineari di due equazioni:
\begin{multicols}{2}
\begin{enumerate}[leftmargin=3.2cm]
    \item
    \[
    \begin{cases}
    \dfrac{x-1}{2}-\dfrac{y+1}{3}=0\\
    \dfrac{x+1}{4}+\dfrac{1}{2}y=6
    \end{cases}
    \]

    \item
    \[
    \begin{cases}
    \dfrac{2x+y}{3}-\dfrac{x-y}{2}=3\\
    \dfrac{x+y}{2}-\dfrac{1}{3}x=2
    \end{cases}
    \]

    \item
    \[
    \begin{cases}
    \dfrac{1}{2}x+\dfrac{2y-1}{3}=2\\
    \dfrac{3x-2}{4}-\dfrac{1}{2}y=0
    \end{cases}
    \]

    \item
    \[
    \begin{cases}
    \dfrac{2(x-1)}{3}+\dfrac{1}{4}y=2\\
    \dfrac{x+y}{3}-\dfrac{y-x}{7}=2
    \end{cases}
    \]
\end{enumerate}
\end{multicols}
\begin{solution}
\begin{enumerate}
    \item $(x,y)=(7,8)$
    \item $(x,y)=(3,3)$
    \item $(x,y)=(2,2)$
    \item $(x,y)=(1,8)$
\end{enumerate}
\end{solution}
\end{question}

%----------------------------------------------------%

%====================================================%
% SEZIONE 3: Problemi risolvibili con i sistemi
%====================================================%

\begin{question}[ID=3.1, subtitle={Somma e differenza di due numeri}, difficulty=facile, tags={sistemi, problemi}]
La somma di due numeri è $48$ e la loro differenza è $12$. Determina i due numeri.
\begin{solution}
Indichiamo con $x$ il numero maggiore e con $y$ il minore:
\[
\begin{cases}
x+y=48\\
x-y=12
\end{cases}
\]
Sommando membro a membro: $2x=60$, quindi $x=30$ e $y=48-30=18$.

I due numeri sono $30$ e $18$.
\end{solution}
\end{question}

%----------------------------------------------------%

\begin{question}[ID=3.2, subtitle={Biglietti del cinema}, difficulty=media, tags={sistemi, problemi}]
Al cinema, $3$ biglietti interi e $2$ ridotti costano in tutto $37$ euro, mentre $1$ biglietto intero e $4$ ridotti costano $29$ euro. Quanto costano un biglietto intero e uno ridotto?
\begin{solution}
Sia $x$ il prezzo del biglietto intero e $y$ quello del ridotto:
\[
\begin{cases}
3x+2y=37\\
x+4y=29
\end{cases}
\]
Dalla seconda: $x=29-4y$. Sostituendo nella prima: $87-12y+2y=37$, cioè $10y=50$, quindi $y=5$ e $x=9$.

Il biglietto intero costa $9$ euro, il ridotto $5$ euro.
\end{solution}
\end{question}

%----------------------------------------------------%

\begin{question}[ID=3.3, subtitle={Dimensioni di un rettangolo}, difficulty=media, tags={sistemi, problemi}]
Un rettangolo ha il perimetro di $56\,\text{cm}$ e la base supera di $4\,\text{cm}$ il doppio dell'altezza. Calcola le dimensioni del rettangolo.
\begin{solution}
Siano $b$ la base e $h$ l'altezza. Il semiperimetro è $28\,\text{cm}$:
\[
\begin{cases}
b+h=28\\
b=2h+4
\end{cases}
\]
Sostituendo: $2h+4+h=28$, cioè $3h=24$, quindi $h=8$ e $b=20$.

La base misura $20\,\text{cm}$ e l'altezza $8\,\text{cm}$.
\end{solution}
\end{question}

%----------------------------------------------------%

\begin{question}[ID=3.4, subtitle={Monete in un salvadanaio}, difficulty=difficile, tags={sistemi, problemi}]
In un salvadanaio ci sono $30$ monete, alcune da $50$ centesimi e altre da $2$ euro, per un totale di $36$ euro. Quante monete ci sono di ciascun tipo?
\begin{solution}
Sia $x$ il numero di monete da $0{,}50$ euro e $y$ quello delle monete da $2$ euro:
\[
\begin{cases}
x+y=30\\
0{,}5x+2y=36
\end{cases}
\]
Dalla prima: $x=30-y$. Sostituendo: $15-0{,}5y+2y=36$, cioè $1{,}5y=21$, quindi $y=14$ e $x=16$.

Ci sono $16$ monete da $50$ centesimi e $14$ monete da $2$ euro.
\end{solution}
\end{question}

%----------------------------------------------------%

\begin{question}[ID=3.5, subtitle={Età di padre e figlio}, difficulty=difficile, tags={sistemi, problemi}]
Un padre ha oggi un'età pari al triplo di quella del figlio. Fra $12$ anni l'età del padre sarà il doppio di quella del figlio. Quanti anni hanno oggi il padre e il figlio?
\begin{solution}
Siano $p$ l'età del padre e $f$ quella del figlio:
\[
\begin{cases}
p=3f\\
p+12=2(f+12)
\end{cases}
\]
Sostituendo: $3f+12=2f+24$, quindi $f=12$ e $p=36$.

Il figlio ha $12$ anni e il padre $36$ anni.
\end{solution}
\end{question}

%----------------------------------------------------%

\begin{question}[ID=3.6, subtitle={Miscela di due soluzioni}, difficulty=difficile, tags={sistemi, problemi, percentuali}]
Un chimico vuole ottenere $40\,\text{litri}$ di una soluzione al $32\%$ mescolando una soluzione al $20\%$ con una soluzione al $50\%$. Quanti litri di ciascuna deve utilizzare?
\begin{solution}
Siano $x$ i litri della soluzione al $20\%$ e $y$ quelli della soluzione al $50\%$. Il soluto totale deve essere il $32\%$ di $40\,\text{litri}$, cioè $12{,}8\,\text{litri}$:
\[
\begin{cases}
x+y=40\\
0{,}2x+0{,}5y=12{,}8
\end{cases}
\]
Dalla prima: $x=40-y$. Sostituendo: $8-0{,}2y+0{,}5y=12{,}8$, cioè $0{,}3y=4{,}8$, quindi $y=16$ e $x=24$.

Servono $24\,\text{litri}$ di soluzione al $20\%$ e $16\,\text{litri}$ di soluzione al $50\%$.
\end{solution}
\end{question}

%----------------------------------------------------%

\begin{question}[ID=3.7, subtitle={Battello e corrente del fiume}, difficulty=difficile, tags={sistemi, problemi, moto}]
Un battello percorre $36\,\text{km}$ controcorrente in $3$ ore e gli stessi $36\,\text{km}$ a favore di corrente in $2$ ore. Calcola la velocità del battello in acqua ferma e la velocità della corrente.
\begin{solution}
Siano $v$ la velocità del battello e $c$ quella della corrente. Controcorrente la velocità effettiva è $v-c=\dfrac{36}{3}=12\,\text{km/h}$, a favore è $v+c=\dfrac{36}{2}=18\,\text{km/h}$:
\[
\begin{cases}
v+c=18\\
v-c=12
\end{cases}
\]
Sommando membro a membro: $2v=30$, quindi $v=15$ e $c=3$.

Il battello va a $15\,\text{km/h}$ in acqua ferma e la corrente scorre a $3\,\text{km/h}$.
\end{solution}
\end{question}

%----------------------------------------------------%

\begin{question}[ID=3.8, subtitle={Numero di due cifre}, difficulty=difficile, tags={sistemi, problemi, numeri}]
Un numero di due cifre ha la somma delle cifre uguale a $11$. Scambiando tra loro le cifre si ottiene un numero che è minore di $27$ rispetto a quello di partenza. Determina il numero.
\begin{solution}
Siano $a$ la cifra delle decine e $b$ quella delle unità. Il numero è $10a+b$, quello con le cifre scambiate è $10b+a$:
\[
\begin{cases}
a+b=11\\
(10a+b)-(10b+a)=27
\end{cases}
\]
La seconda equazione diventa $9a-9b=27$, cioè $a-b=3$. Sommando con la prima: $2a=14$, quindi $a=7$ e $b=4$.

Il numero è $74$ (infatti $74-47=27$).
\end{solution}
\end{question}

%----------------------------------------------------%

\begin{question}[ID=3.9, subtitle={Due ciclisti}, difficulty=difficile, tags={sistemi, problemi, moto}]
Due ciclisti partono contemporaneamente da due città distanti $90\,\text{km}$. Se vanno l'uno incontro all'altro, si incontrano dopo $2$ ore. Se invece vanno nello stesso verso, il più veloce raggiunge l'altro dopo $6$ ore. Calcola le velocità dei due ciclisti.
\begin{solution}
Siano $x$ la velocità del ciclista più veloce e $y$ quella dell'altro, in $\text{km/h}$. Nel primo caso le velocità si sommano, nel secondo si sottraggono:
\[
\begin{cases}
2(x+y)=90\\
6(x-y)=90
\end{cases}
\qquad\Longrightarrow\qquad
\begin{cases}
x+y=45\\
x-y=15
\end{cases}
\]
Sommando membro a membro: $2x=60$, quindi $x=30$ e $y=15$.

Le velocità sono $30\,\text{km/h}$ e $15\,\text{km/h}$.
\end{solution}
\end{question}

%----------------------------------------------------%

\begin{question}[ID=3.10, subtitle={Prezzi scontati}, difficulty=difficile, tags={sistemi, problemi, percentuali}]
Due articoli costano complessivamente $120$ euro. Il primo viene scontato del $20\%$ e il secondo del $10\%$, e in tutto si pagano $100$ euro. Quanto costava ciascun articolo prima dello sconto?
\begin{solution}
Siano $x$ e $y$ i prezzi originari. Dopo lo sconto si paga l'$80\%$ di $x$ e il $90\%$ di $y$:
\[
\begin{cases}
x+y=120\\
0{,}8x+0{,}9y=100
\end{cases}
\]
Dalla prima: $x=120-y$. Sostituendo: $96-0{,}8y+0{,}9y=100$, cioè $0{,}1y=4$, quindi $y=40$ e $x=80$.

Il primo articolo costava $80$ euro, il secondo $40$ euro.
\end{solution}
\end{question}

%====================================================%
% SEZIONE 4: Interpretazione grafica e break-even point
%====================================================%

\begin{question}[ID=4.1, subtitle={Break-even di una produzione di magliette}, difficulty=media, tags={sistemi, break-even, grafici}]
Un'azienda produce magliette sostenendo costi fissi di $600$ euro e un costo di $4$ euro per ogni maglietta prodotta. Ogni maglietta viene venduta a $10$ euro.
\begin{enumerate}[label=\alph*),leftmargin=1.8cm]
    \item Scrivi la funzione dei costi $C(x)$ e quella dei ricavi $R(x)$, dove $x$ è il numero di magliette prodotte e vendute.
    \item Determina il break-even point risolvendo il sistema tra le due funzioni.
    \item Rappresenta graficamente $C(x)$ e $R(x)$ nello stesso piano cartesiano e indica per quali valori di $x$ l'azienda è in utile e per quali è in perdita.
\end{enumerate}
\begin{solution}
\begin{enumerate}[label=\alph*),leftmargin=1.8cm]
    \item $C(x)=600+4x$ e $R(x)=10x$.
    \item Il break-even point si ottiene quando costi e ricavi coincidono:
    \[
    \begin{cases}
    y=600+4x\\
    y=10x
    \end{cases}
    \quad\Longrightarrow\quad 10x=600+4x \quad\Longrightarrow\quad 6x=600 \quad\Longrightarrow\quad x=100.
    \]
    Quindi $y=1000$: il break-even point è $B(100;\,1000)$.
    \item Le due rette si disegnano con due punti ciascuna: $C$ passa per $(0;\,600)$ e $(100;\,1000)$, $R$ passa per $(0;\,0)$ e $(100;\,1000)$. Si intersecano in $B(100;\,1000)$.

    Per $x<100$ la retta dei costi sta sopra quella dei ricavi: l'azienda è in \textbf{perdita}. Per $x>100$ la retta dei ricavi sta sopra: l'azienda è in \textbf{utile}. Per $x=100$ non c'è né utile né perdita.
\begin{center}
\begin{tikzpicture}
\begin{axis}[
    width=10cm,
    height=6cm,
    xmin=0, xmax=160,
    ymin=0, ymax=1700,
    axis lines=left,
    xlabel={$x$},
    ylabel={$y$ (euro)},
    grid=both,
]
\addplot[red,thick,domain=0:160] {600+4*x};
\node[red] at (axis cs:140,1230) {C};
\addplot[blue,thick,domain=0:160] {10*x};
\node[blue] at (axis cs:140,1560) {R};
\node[black] at (axis cs:100,1000) {● B (100; 1000)};
\end{axis}
\end{tikzpicture}
\end{center}
\end{enumerate}
\end{solution}
\end{question}

%----------------------------------------------------%

\begin{question}[ID=4.2, subtitle={Break-even di un carretto dei gelati}, difficulty=media, tags={sistemi, break-even, grafici}]
Un gelataio affitta un carretto pagando $450$ euro al mese. Ogni gelato gli costa $0{,}50$ euro di ingredienti e viene venduto a $2{,}50$ euro.
\begin{enumerate}[label=\alph*),leftmargin=1.8cm]
    \item Scrivi le funzioni dei costi e dei ricavi mensili in funzione del numero $x$ di gelati venduti.
    \item Quanti gelati deve vendere al mese per raggiungere il break-even point? A quanto ammontano in quel punto ricavi e costi?
    \item Interpreta graficamente la situazione: come si dispongono le due rette prima e dopo il break-even point?
    \item Scrivi la funzione utile $U(x)=R(x)-C(x)$ e calcola l'utile se vende $300$ gelati.
\end{enumerate}
\begin{solution}
\begin{enumerate}[label=\alph*),leftmargin=1.8cm]
    \item $C(x)=450+0{,}5x$ e $R(x)=2{,}5x$.
    \item Risolviamo il sistema:
    \[
    \begin{cases}
    y=450+0{,}5x\\
    y=2{,}5x
    \end{cases}
    \quad\Longrightarrow\quad 2x=450 \quad\Longrightarrow\quad x=225.
    \]
    Quindi $y=2{,}5\cdot 225=562{,}5$. Servono $225$ gelati al mese e in quel punto costi e ricavi valgono $562{,}5$ euro: $B(225;\,562{,}5)$.
    \item La retta dei costi parte da $(0;\,450)$ ed è poco inclinata (coefficiente angolare $0{,}5$), quella dei ricavi parte dall'origine ed è più inclinata (coefficiente angolare $2{,}5$). Le due rette si intersecano in $B$: a sinistra di $B$ i costi superano i ricavi (perdita), a destra i ricavi superano i costi (utile).
\begin{center}
\begin{tikzpicture}
\begin{axis}[
    width=10cm,
    height=6cm,
    xmin=0, xmax=400,
    ymin=0, ymax=1100,
    axis lines=left,
    xlabel={$x$},
    ylabel={$y$ (euro)},
    grid=both,
]
\addplot[red,thick,domain=0:400] {450+0.5*x};
\node[red] at (axis cs:350,560) {C};
\addplot[blue,thick,domain=0:400] {2.5*x};
\node[blue] at (axis cs:350,960) {R};
\node[black] at (axis cs:225,562.5) {● B (225; 562,5)};
\end{axis}
\end{tikzpicture}
\end{center}
    \item $U(x)=2{,}5x-(450+0{,}5x)=2x-450$. Per $x=300$: $U(300)=600-450=150$ euro.
\end{enumerate}
\end{solution}
\end{question}

%----------------------------------------------------%

\begin{question}[ID=4.3, subtitle={Confronto tra due piani telefonici}, difficulty=difficile, tags={sistemi, break-even, grafici}]
Un operatore telefonico propone due tariffe mensili. Il piano A costa $5$ euro di canone fisso più $0{,}10$ euro al minuto. Il piano B costa $20$ euro di canone fisso più $0{,}04$ euro al minuto.
\begin{enumerate}[label=\alph*),leftmargin=1.8cm]
    \item Scrivi la spesa mensile $S_A(x)$ e $S_B(x)$ in funzione dei minuti di conversazione $x$.
    \item Determina per quanti minuti le due tariffe costano uguale e qual è in quel caso la spesa.
    \item Rappresenta graficamente le due rette e, osservando il grafico, stabilisci quale piano conviene a seconda dei minuti utilizzati.
\end{enumerate}
\begin{solution}
\begin{enumerate}[label=\alph*),leftmargin=1.8cm]
    \item $S_A(x)=5+0{,}10x$ e $S_B(x)=20+0{,}04x$.
    \item Risolviamo il sistema:
    \[
    \begin{cases}
    y=5+0{,}10x\\
    y=20+0{,}04x
    \end{cases}
    \quad\Longrightarrow\quad 0{,}06x=15 \quad\Longrightarrow\quad x=250.
    \]
    Quindi $y=5+0{,}10\cdot 250=30$. Le tariffe sono uguali a $250$ minuti, con una spesa di $30$ euro: punto $P(250;\,30)$.
    \item La retta A parte più in basso (ordinata all'origine $5$) ma è più inclinata (coefficiente angolare $0{,}10$ contro $0{,}04$), quindi dopo l'intersezione si trova sopra la retta B.

    Per $x<250$ la retta A è sotto B: conviene il \textbf{piano A}. Per $x>250$ la retta B è sotto A: conviene il \textbf{piano B}. Per $x=250$ i due piani sono equivalenti.
\begin{center}
\begin{tikzpicture}
\begin{axis}[
    width=10cm,
    height=6cm,
    xmin=0, xmax=500,
    ymin=0, ymax=62,
    axis lines=left,
    xlabel={$x$ (minuti)},
    ylabel={$y$ (euro)},
    grid=both,
]
\addplot[red,thick,domain=0:500] {5+0.1*x};
\node[red] at (axis cs:450,51) {A};
\addplot[blue,thick,domain=0:500] {20+0.04*x};
\node[blue] at (axis cs:450,41) {B};
\node[black] at (axis cs:250,30) {● P (250; 30)};
\end{axis}
\end{tikzpicture}
\end{center}
\end{enumerate}
\end{solution}
\end{question}

%----------------------------------------------------%

\begin{question}[ID=4.4, subtitle={Noleggio auto: quale agenzia conviene}, difficulty=difficile, tags={sistemi, break-even, grafici}]
Due agenzie noleggiano la stessa auto. L'agenzia Alfa chiede $50$ euro di quota fissa più $0{,}20$ euro per ogni chilometro percorso. L'agenzia Beta chiede $20$ euro di quota fissa più $0{,}30$ euro per ogni chilometro.
\begin{enumerate}[label=\alph*),leftmargin=1.8cm]
    \item Scrivi le due funzioni di costo in funzione dei chilometri $x$ percorsi.
    \item Dopo quanti chilometri i due costi coincidono? Quanto si spende in quel caso?
    \item Traccia i grafici e spiega, usando il punto di intersezione, per quali percorrenze conviene ciascuna agenzia.
    \item Un cliente deve percorrere $450\,\text{km}$. Quale agenzia sceglie e quanto risparmia rispetto all'altra?
\end{enumerate}
\begin{solution}
\begin{enumerate}[label=\alph*),leftmargin=1.8cm]
    \item $C_A(x)=50+0{,}2x$ e $C_B(x)=20+0{,}3x$.
    \item Risolviamo il sistema:
    \[
    \begin{cases}
    y=50+0{,}2x\\
    y=20+0{,}3x
    \end{cases}
    \quad\Longrightarrow\quad 0{,}1x=30 \quad\Longrightarrow\quad x=300.
    \]
    Quindi $y=50+0{,}2\cdot 300=110$. I costi coincidono a $300\,\text{km}$ e valgono $110$ euro: punto $P(300;\,110)$.
    \item Le rette partono da $(0;\,50)$ e $(0;\,20)$. Beta parte più in basso ma cresce più rapidamente. Per $x<300$ la retta di Beta è sotto quella di Alfa: conviene \textbf{Beta}. Per $x>300$ la retta di Alfa è sotto: conviene \textbf{Alfa}.
\begin{center}
\begin{tikzpicture}
\begin{axis}[
    width=10cm,
    height=6cm,
    xmin=0, xmax=500,
    ymin=0, ymax=190,
    axis lines=left,
    xlabel={$x$ (km)},
    ylabel={$y$ (euro)},
    grid=both,
]
\addplot[red,thick,domain=0:500] {50+0.2*x};
\node[red] at (axis cs:440,130) {Alfa};
\addplot[blue,thick,domain=0:500] {20+0.3*x};
\node[blue] at (axis cs:440,165) {Beta};
\node[black] at (axis cs:300,110) {● P (300; 110)};
\end{axis}
\end{tikzpicture}
\end{center}
    \item Poiché $450>300$ conviene Alfa. Costi: $C_A(450)=50+90=140$ euro, $C_B(450)=20+135=155$ euro. Il cliente risparmia $15$ euro.
\end{enumerate}
\end{solution}
\end{question}

%----------------------------------------------------%

\begin{question}[ID=4.5, subtitle={Break-even e utile desiderato}, difficulty=difficile, tags={sistemi, break-even, grafici}]
Un laboratorio artigianale produce zaini. Ogni mese sostiene costi fissi di $1\,200$ euro e un costo di produzione di $8$ euro per zaino. Ogni zaino viene venduto a $14$ euro.
\begin{enumerate}[label=\alph*),leftmargin=1.8cm]
    \item Determina il break-even point e rappresentalo graficamente come intersezione delle rette dei costi e dei ricavi.
    \item Quanti zaini deve vendere il laboratorio per ottenere un utile mensile di $900$ euro?
    \item Il prezzo di vendita sale a $16$ euro. Come cambia il break-even point? Come cambia la retta dei ricavi nel grafico?
\end{enumerate}
\begin{solution}
\begin{enumerate}[label=\alph*),leftmargin=1.8cm]
    \item Le funzioni sono $C(x)=1200+8x$ e $R(x)=14x$:
    \[
    \begin{cases}
    y=1200+8x\\
    y=14x
    \end{cases}
    \quad\Longrightarrow\quad 6x=1200 \quad\Longrightarrow\quad x=200.
    \]
    Quindi $y=2800$: break-even point $B(200;\,2800)$. Nel grafico la retta dei costi parte da $(0;\,1200)$, quella dei ricavi dall'origine, e si incontrano in $B$.
\begin{center}
\begin{tikzpicture}
\begin{axis}[
    width=10cm,
    height=6cm,
    xmin=0, xmax=300,
    ymin=0, ymax=4500,
    axis lines=left,
    xlabel={$x$},
    ylabel={$y$ (euro)},
    grid=both,
]
\addplot[red,thick,domain=0:300] {1200+8*x};
\node[red] at (axis cs:260,3050) {C};
\addplot[blue,thick,domain=0:300] {14*x};
\node[blue] at (axis cs:260,3800) {R};
\node[black] at (axis cs:200,2800) {● B (200; 2800)};
\end{axis}
\end{tikzpicture}
\end{center}
    \item L'utile è $U(x)=R(x)-C(x)=6x-1200$. Imponendo $U(x)=900$: $6x=2100$, quindi $x=350$. Il laboratorio deve vendere $350$ zaini.

    Graficamente, si cerca il punto in cui la retta dei ricavi supera quella dei costi di esattamente $900$ euro.
    \item Con $R(x)=16x$ si ha $16x=1200+8x$, cioè $8x=1200$, quindi $x=150$ e $y=2400$: il nuovo break-even point è $B'(150;\,2400)$. Il break-even viene raggiunto con meno zaini (da $200$ a $150$). Nel grafico la retta dei ricavi diventa più inclinata (coefficiente angolare da $14$ a $16$) e interseca prima la retta dei costi, che resta invariata.
\end{enumerate}
\end{solution}
\end{question}
